%=============================================================================
%  AD_COURSEWORK_GROUP_X.tex
%
%  Advanced Derivatives -- Coursework report template (2026-27)
%
%  How to use it:
%   1. Replace "X" in the group number below and fill in the names/IDs.
%   2. Work through Questions 1-4. Every grey box is a placeholder for a
%      figure, a table or a code listing. The italic line under each heading
%      restates what the question asks -- delete it before you submit.
%   3. Delete the appendix "Hints" section from your final PDF.
%
%  The report (excluding the cover page) must be at most 10 pages.
%=============================================================================
\documentclass[12pt,a4paper]{article}

\usepackage[T1]{fontenc}
\usepackage[utf8]{inputenc}
\usepackage[a4paper,margin=2.5cm]{geometry}
\usepackage{amsmath,amssymb}
\usepackage{graphicx}
\usepackage{xcolor}
\usepackage{xspace}
\usepackage{booktabs}
\usepackage{tabularx}
\usepackage{float}
\usepackage{caption}
\usepackage{listings}
\usepackage{enumitem}
\usepackage{hyperref}

\definecolor{linkblue}{RGB}{0,70,140}
\hypersetup{colorlinks=true, linkcolor=linkblue, urlcolor=linkblue,
            citecolor=linkblue}

\graphicspath{{figures/}}

\lstset{
  language=Python,
  basicstyle=\ttfamily\small,
  keywordstyle=\color{blue},
  commentstyle=\color{green!50!black},
  stringstyle=\color{red!70!black},
  showstringspaces=false,
  numbers=left,
  numberstyle=\tiny,
  stepnumber=1,
  frame=single,
  breaklines=true
}

%-----------------------------------------------------------------------------
%  Template macros -- replace the boxes, do not redefine them
%-----------------------------------------------------------------------------

% Framed "Answer." line that opens every sub-question.
\newcommand{\subq}[2]{%
  \subsection*{#1}%
  \noindent\textit{#2}\par\medskip
  \noindent\textbf{Answer.}\par}

% Placeholder for a figure.  Usage: \figplaceholder[width]{caption}{label}
\newcommand{\figplaceholder}[3][0.85\textwidth]{%
  \begin{figure}[H]
    \centering
    \fbox{\parbox[c][3.4cm][c]{#1}{\centering\color{gray}\small
      FIGURE PLACEHOLDER\\[2mm]
      replace this box with\\[1mm]
      \texttt{\textbackslash includegraphics[width=\textbackslash linewidth]\{...\}}}}
    \caption{#2}
    \label{#3}
  \end{figure}}

% Placeholder for a table.  Usage: \tabplaceholder{caption}{label}
\newcommand{\tabplaceholder}[2]{%
  \begin{table}[H]
    \centering
    \fbox{\parbox[c][2.4cm][c]{0.85\textwidth}{\centering\color{gray}\small
      TABLE PLACEHOLDER\\[2mm]
      paste the output of \texttt{df.to\_latex(...)}}}
    \caption{#1}
    \label{#2}
  \end{table}}

% Placeholder for a code listing.  Usage: \codeslot{description}
\newcommand{\codeslot}[1]{%
  \begin{center}
    \fbox{\parbox{0.9\textwidth}{\centering\color{gray}\small
      CODE PLACEHOLDER: #1\\[1mm]
      use \texttt{\textbackslash lstinputlisting[language=Python]\{...\}}}}
  \end{center}}

%-----------------------------------------------------------------------------
\begin{document}

%-----------------------------------------------------------------------------
%  COVER PAGE
%-----------------------------------------------------------------------------
\begin{titlepage}
\centering
\vspace*{1.5cm}

{\Large\bfseries EDHEC Business School}\\[3mm]
{\large M.Sc. in Financial Engineering}\\[18mm]

{\LARGE\bfseries Advanced Derivatives}\\[5mm]
{\LARGE Coursework Report}\\[6mm]
{\Large Academic Year 2026--27}\\[18mm]

{\Large\bfseries Group X}\\[10mm]

\begin{tabular}{@{}cll@{}}
\toprule
\# & Name & Student ID\\
\midrule
1 & \hspace{7cm} & \hspace{3cm}\\
2 & \hspace{7cm} & \hspace{3cm}\\
3 & \hspace{7cm} & \hspace{3cm}\\
\bottomrule
\end{tabular}

\vfill
{\large Deadline: Tuesday 1 December 2026}\\[3mm]
{\normalsize Upload as \texttt{AD\_COURSEWORK\_GROUP\_X.zip}}
\vspace*{1.5cm}
\end{titlepage}

%-----------------------------------------------------------------------------
%  SUBMISSION CHECKLIST
%-----------------------------------------------------------------------------
\section*{Submission checklist}
\begin{center}
\begin{itemize}[leftmargin=*,itemsep=2pt]
  \item The ZIP contains this PDF report and \textbf{one} notebook or
        \texttt{.py} file with all the code for Questions 1--4.
  \item The report answers all questions \textbf{in order}, with every table
        and figure printed in the PDF (do not refer to the notebook for them).
  \item Excluding this cover page, the report is at most \textbf{12 pages}.
  \item The code is split clearly by question, commented, and your results are
        explained in plain English.
  \item No Word documents. PDF only.
\end{itemize}
\end{center}

\section*{Explanation \LaTeX}
\begin{itemize}
  \item You can use this template with a local \LaTeX\xspace installation (e.g.
    TexMaker (editor) and a local \LaTeX\xspace installation (TexLive, MikTeX,
    MacTeX)) or usign \href{overleaf.com}{Overleaf} (easiest).
  \item For those that would like to learn more See
    \href{https://www.youtube.com/watch?v=0ivLZh9xK1Q&list=PL1D4EAB31D3EBC449}{LaTeX
    course} for a good course.
  \item Feel free to use an AI assistant to help you with the \LaTeX\xspace part.
\end{itemize}

\newpage
%=============================================================================
\section*{Question 1: Black-Scholes-Merton hedging}
%=============================================================================

\subq{1(a) Using the FinancePy library}
{Install FinancePy and call the model functions
(\texttt{bs\_value}, \texttt{bs\_delta}, \ldots) directly instead of going
through the \texttt{EquityOptions} wrapper, which is much slower.}
\codeslot{the import and the direct library calls you use}

\subq{1(b) The \texttt{DeltaHedge} function}
{Write a function that simulates delta hedging of a European put from trade
date to expiry using a self-financing portfolio. Inputs: $K$, $S$, $r$,
$\mu$, $\sigma$, $T$ and the number of hedges per year $N$. Output: a tuple
with the terminal stock price, the option payoff, the realised variance of
returns, the replicating error (portfolio value minus payoff) and the
average gamma over the hedging period.}
\codeslot{the full definition of \texttt{DeltaHedge}}

\subq{1(c) Running 10,000 paths}
{Write a second function that calls \texttt{DeltaHedge} over 10,000
independent paths and returns the results in a convenient (e.g.\ DataFrame)
form. Include a clear listing of the code.}
\codeslot{the path-loop / vectorised wrapper}

\subq{1(d) Hedging error versus terminal stock price}
{For $S(0)=100$, $K=100$, $r=4\%$, $T=1$, $\sigma=20\%$ and $\mu=8\%$,
scatter the hedging error ($y$) against the terminal stock price ($x$) for
$N=12$, $N=52$ and $N=252$, using different colours or symbols.}
\figplaceholder{Hedging error against terminal stock price for $N=12$, $52$
and $252$. Caption must state the parameters.}{fig:q1-error}

\subq{1(e) Mean and variance of the hedging error}
{For each $N$, report the mean and the variance of the hedging error over
10,000 paths in a simple table.}
\tabplaceholder{Mean and variance of the hedging error over 10,000 paths.}
{tab:q1-moments}

\subq{1(f) Realised volatility and the replication error}
{Scatter the realised volatility against the replication error and explain
the pattern you see.}
\figplaceholder{Replication error against realised volatility.}{fig:q1-realvol}

\subq{1(g) Average gamma and the replication error}
{In addition to the delta, compute the gamma at each hedging date. Scatter
the replication error against the average gamma along each path and explain
the pattern.}
\figplaceholder{Replication error against average gamma.}{fig:q1-gamma}

\subq{1(h) Drift sensitivity}
{Repeat the exercise for $\mu=2.5\%,\,5.0\%,\,7.5\%,\,10\%$ with $N=52$ and
10,000 paths. Report the mean absolute error and the variance of the
hedging error in a table. Is the effect of the drift small or large?}
\tabplaceholder{Hedging error statistics for different drift assumptions
($N=52$, 10,000 paths).}{tab:q1-drift}

\newpage
%=============================================================================
\section*{Question 2: Time-varying volatility}
%=============================================================================

\subq{2(a) A two-period volatility term structure}
{Value a put with $S(0)=100$, $K=100$, $r=4\%$ and $T=1$ when the
volatility is $4\%$ over the first half of the option's life and $28\%$ over
the second half. Compare with the flat $\sigma=20\%$ value and explain why
they agree.}
\tabplaceholder{Put values: piecewise volatility ($4\%$ then $28\%$) versus
flat $20\%$.}{tab:q2-termstructure}

\subq{2(b) \texttt{DeltaHedgeTV}}
{Adapt \texttt{DeltaHedge} so that the volatility can vary over time
(\texttt{DeltaHedgeTV}). Explain how the volatility used for pricing and for
the stock path changes as time to maturity decreases.}
\codeslot{the full definition of \texttt{DeltaHedgeTV}}

\subq{2(c) Hedging performance with time-varying volatility}
{Repeat the hedging experiment of 1(d) for the piecewise volatility case and
compare with the constant volatility results.}
\figplaceholder{Hedging error against terminal stock price, piecewise versus
constant volatility.}{fig:q2-hedge}
\tabplaceholder{Mean and variance of the hedging error: constant versus
time-varying volatility.}{tab:q2-compare}

\subq{2(d) Reversing the volatility structure}
{Now let the first half have $\sigma=28\%$ and the second half $\sigma=4\%$.
Repeat the analysis and report the results.}
\tabplaceholder{Hedging error statistics when the high-volatility period
comes first.}{tab:q2-reversed}

\subq{2(e) Comparing the three cases}
{Explain the differences (if any) between constant volatility, low volatility
first and low volatility second. Support the discussion with your tables and
figures.}

\newpage
%=============================================================================
\section*{Question 3: Asian options and the control variate technique}
%=============================================================================

\subq{3(a) Pricing an arithmetic average Asian call}
{Write a Monte-Carlo function that prices a 6-month arithmetic average Asian
call. Averaging starts after 3 months and runs daily (use 90 days) to
expiry. Include the code in the PDF.}
\codeslot{the arithmetic Asian Monte-Carlo pricing function}

\subq{3(b) A first price}
{Price the option for $S(0)=100$, $K=100$, $r=4\%$, $\sigma=25\%$ with
$P=10{,}000$ paths. Report the price and its standard error.}

\subq{3(c) Adding the geometric average}
{Extend the function to price the geometric average option \emph{on the same
paths and with the same random numbers}, returning both prices.}
\codeslot{the joint arithmetic/geometric pricing function}

\subq{3(d) Variances and covariance}
{Extend the function so that it returns five numbers: the geometric price,
the arithmetic price, their two variances and the covariance of the
discounted payoffs across the $P$ paths. Provide a code listing.}
\codeslot{the function returning prices, variances and covariance}

\subq{3(e) Numerical results}
{Compute the arithmetic and geometric prices for $S(0)=100$, $r=5\%$,
$q=1\%$, $\sigma=30\%$, $K=100$ and $P=5{,}000$ paths.}
\tabplaceholder{Arithmetic and geometric Asian prices, variances and
covariance ($P=5{,}000$).}{tab:q3-results}

\subq{3(f) The control variate estimate}
{With $T_0=0$, value the geometric option analytically and use it as a
control variate for the arithmetic price. State the estimator you use and
report the best estimate.}

\subq{3(g) Paths needed with the control variate}
{Estimate how many paths the control variate method needs to reach the same
variance (or confidence interval width) as the plain arithmetic simulation
with $P=10{,}000$ paths. Show the calculation.}
\tabplaceholder{Variance of the plain and control-variate estimators and the
equivalent number of paths.}{tab:q3-paths}

\newpage
%=============================================================================
\section*{Question 4: The delta of the Asian option}
%=============================================================================

\subq{4(a) Finite-difference delta}
{Compute the delta of the Asian option by finite differences,
\begin{equation*}
  \frac{\partial f}{\partial S(0)} \approx
  \frac{f(S(0)+h) - f(S(0))}{h},
\end{equation*}
plot the delta for several bump sizes $h$ and explain which random numbers
you reuse.}
\figplaceholder{Finite-difference delta of the Asian option as a function of
the bump size $h$.}{fig:q4-fd}
\tabplaceholder{Delta estimates for the bump sizes considered.}{tab:q4-fd}

\subq{4(b) Pathwise delta}
{Estimate the delta with the pathwise estimator
\begin{equation*}
  \frac{\partial f}{\partial S(0)}
  = \mathbb{E}\left[e^{-rT}\,\mathbb{I}(\bar{S}>K)\,
    \frac{\bar{S}}{S(0)}\right],
\end{equation*}
using the sample average over the Monte-Carlo paths and \emph{without}
control variates.}
\codeslot{the pathwise delta estimator}
\tabplaceholder{Pathwise delta estimate and its standard error.}{tab:q4-pathwise}

\subq{4(c) Control variate for the delta}
{Investigate whether the geometric-average control variate also helps for
the delta. Is it as effective as it is for the option price? Quantify your
answer.}
\tabplaceholder{Variance reduction for the delta: plain versus control
variate estimator.}{tab:q4-cv}

%=============================================================================
\newpage
\section*{Appendix: Hints for producing figures and tables from Python}
%=============================================================================
Delete this appendix before submitting. It only shows how to get the output
of your notebook into this document.

\subsection*{1. Figures}

Save the figure from the notebook \emph{before} you call
\texttt{plt.show()}, into a \texttt{figures/} folder next to this file:

\begin{lstlisting}
import matplotlib.pyplot as plt

fig, ax = plt.subplots(figsize=(6, 4))
ax.scatter(S_T, error, s=4)
ax.set_xlabel("Terminal stock price")
ax.set_ylabel("Hedging error")

# PDF keeps the plot sharp at any zoom; use dpi=300 with .png if you prefer
fig.savefig("figures/q1_error.pdf", bbox_inches="tight")
plt.show()
\end{lstlisting}

Then replace a placeholder box with:

\begin{lstlisting}[language=TeX]
\begin{figure}[H]
  \centering
  \includegraphics[width=0.8\linewidth]{q1_error.pdf}
  \caption{Hedging error against terminal stock price.}
  \label{fig:q1-error}
\end{figure}
\end{lstlisting}

\texttt{\textbackslash graphicspath\{\{figures/\}\}} is already set, so the
file name alone is enough. Refer to the figure in the text with
\texttt{\textbackslash ref\{fig:q1-error\}}.

\subsection*{2. Tables from a pandas DataFrame}

Round and rename first, then let pandas write the LaTeX:

\begin{lstlisting}
summary = df.groupby("N")["error"].agg(["mean", "var"]).round(4)
summary.columns = ["Mean", "Variance"]

summary.to_latex("tables/q1_summary.tex",
                 index=False,
                 float_format="%.4f",
                 caption="Hedging error over 10,000 paths.",
                 label="tab:q1-moments")
\end{lstlisting}

Two things to know:
\begin{itemize}[leftmargin=*,itemsep=2pt]
  \item Passing \texttt{caption=} and \texttt{label=} makes pandas emit a
        complete \texttt{table} float, so in LaTeX you only write
        \texttt{\textbackslash input\{tables/q1\_summary.tex\}} at top level.
  \item Without \texttt{caption=} you get only the \texttt{tabular}, which you
        must wrap yourself (and then paste the output instead of using
        \texttt{\textbackslash input} if you want a single self-contained
        file).
\end{itemize}

\begin{lstlisting}[language=TeX]
\begin{table}[H]
  \centering
  \input{tables/q1_summary.tex}
\end{table}
\end{lstlisting}

\texttt{booktabs} rules (\texttt{\textbackslash toprule},
\texttt{\textbackslash midrule}, \texttt{\textbackslash bottomrule}) are
already loaded; use them if you build a table by hand. Keep
\texttt{float\_format="\%.4f"} so every column has the same precision, and
\texttt{index=False} unless the index is meaningful.

\subsection*{3. Code listings}

Do not retype code. Point LaTeX at the submitted file and, if needed, select
the relevant lines:

\begin{lstlisting}[language=TeX]
\lstinputlisting[language=Python, firstline=42, lastline=78]{code.py}
\end{lstlisting}

This guarantees the PDF, the notebook and the \texttt{.py} file never
disagree.

\subsection*{4. Formatting reminders}
\begin{itemize}[leftmargin=*,itemsep=2pt]
  \item Use \texttt{[H]} (from the \texttt{float} package) to keep a float
        where you put it, and \texttt{\textbackslash clearpage} if a large
        figure or table leaves a bad page break.
  \item Number tables and figures consistently and refer to them by
        \texttt{\textbackslash ref}; never write ``the table below''.
  \item Ten pages is a hard limit: the grey boxes here are guides, not
        content.
\end{itemize}

\end{document}
